\section{Evaluation Plan: ModuleComposeBench}

ModuleComposeBench evaluates whether a model can integrate existing RTL/IP modules, not merely generate standalone Verilog. Tasks include wrapper generation, stream chaining, bus attachment, width adaptation, protocol bridging, CDC insertion, and unsafe ambiguity rejection.

Metrics include first-pass compose success, repair turns, lint pass, simulation pass, formal pass, unsafe rejection rate, QoR delta, and connection entropy reduction. Baselines include direct Verilog prompting, SystemVerilog-interface prompting, MICO AST prompting without compiler-feedback repair, and MICO AST prompting with compiler-feedback repair.

\begin{table}[t]
  \centering
  \caption{Evaluation dimensions}
  \begin{tabular}{lll}
    \toprule
    Dimension & Direct RTL & MICO \\
    \midrule
    Output & RTL text & Typed compose graph \\
    Error boundary & Simulation/review & Compiler + RTL tools \\
    Repair unit & Code rewrite & JSON AST patch \\
    CDC handling & Often implicit & Explicit adapter \\
    \bottomrule
  \end{tabular}
\end{table}

\subsection{Benchmark Protocol}

Each benchmark task contains a module inventory, protocol/interface library, optional adapter library, natural-language integration request, expected safety constraints, and reference validation collateral. A run records the model profile, prompt hash, first-pass compiler result, repair turns, lint result, simulation/formal result when available, unsafe rejection result for negative tasks, and generated artifacts. The same task is evaluated across baselines so that improvements can be attributed to the representation and repair loop rather than task selection.

\begin{table}[t]
  \centering
  \caption{ModuleComposeBench task strata}
  \begin{tabular}{lll}
    \toprule
    Level & Focus & Example \\
    \midrule
    L1 & direct streams & FIFO chain \\
    L2 & widths/params & width adapter \\
    L3 & latency/protocol seeds & skid/pipeline checks \\
    L4 & CDC/RDC & async FIFO \\
    L5 & bus/register seeds & control/status wrapper \\
    L6 & subsystem seeds & multi-block chain \\
    \bottomrule
  \end{tabular}
\end{table}

\subsection{Current Seed Results}

The current artifact includes 57 deterministic seed tasks: 31 positive composition tasks and 26 negative rejection tasks. The manifest spans L1--L6, including same-domain stream wiring, width adaptation, latency/backpressure seeds, CDC/RDC adapter tasks, bus/register wrapper seeds, subsystem seeds, and unsafe examples for direct CDC, width mismatches, reversed direction, missing or invalid contract guarantees, unknown declarations, duplicate declarations, protocol mismatches, and adapter misuse. The L3/L5/L6 entries are still seed approximations over the committed smoke RTL collateral rather than dedicated subsystem case studies, but they validate that the compiler, JSON AST path, ready/valid v0 contract checks, SystemVerilog/SVA emitters, traceability report, Verilator lint, Icarus elaboration, Yosys elaboration, positive-seed Icarus simulations, selected bounded SymbiYosys proofs, structural QoR aggregation, and benchmark aggregation path are reproducible.

\begin{table}[t]
  \centering
  \caption{Seed artifact smoke results}
  \begin{tabular}{lll}
    \toprule
    Metric & Result & Scope \\
    \midrule
    Expected outcome & 57/57 & all seed tasks \\
    Compose pass & 31/31 & positive tasks \\
    Lint/elab pass & 31/31 & positive tasks \\
    Simulation pass & 4/4 & harness-enabled positives \\
    Selected formal pass & 2/2 & direct + width seeds \\
    Structural QoR & 4/4 & reference-enabled positives \\
    Unsafe rejection & 26/26 & negative tasks \\
    JSON AST path & 57/57 & source-to-AST checks \\
    \bottomrule
  \end{tabular}
\end{table}

The paper therefore does not claim final LLM pass-rate improvements yet. The CCF-A submission target still requires running multiple LLM baselines, extending simulation and formal harnesses, adding dedicated subsystem case studies, and reporting richer QoR deltas for adapter-heavy tasks. The value of the current result is reproducibility: a clean checkout can run the Rust tests, emit wrappers and SVA skeletons, reject unsafe negative examples, execute supported positive-seed simulations, prove the selected direct and width seed properties, aggregate machine-readable benchmark JSON including structural QoR CSV/TeX summaries, and pass open-source EDA smoke checks in a controlled Docker environment.
